% Abhakseite „Das kann ich“ – Zone nur im Gesamt (\mitzone)
%% TODO Ich-kann-Sätze: Platzhalter wie in den Aufgabendateien
\begin{abhakseite}
\ifdefined\mitzone
\abhakgruppe{Kennst du schon}
\abhak{1}{Bestimmte Integrale berechnen (Stammfunktion, Hauptsatz, auch mit Rechner)}
\abhak{2}{Die Rate als Ableitung lesen}
\abhak{3}{Vorzeichen einer Funktion über Nullstellen bestimmen}
\abhak{4}{Einheiten umrechnen}
\abhak{5}{Rate mal Zeit als Menge deuten (je-desto, proportionale Grundvorstellung)}
\abhak{6}{Einheiten umrechnen – Fehler finden}
\abhak{7}{Einheiten umrechnen – selbst rechnen}
\fi
\abhakgruppe{Einheit 1 · Bestand aus Rate}
\abhak{8}{Bestand aus Rate}
\abhak{9}{Zeitpunkt mit gleichem Bestand wie zu Beginn über das Integral der Änderungsrate gleich null untersuchen}
\abhak{10}{Bestand aus Rate – Fehler finden, Begründen, Darstellungswechsel, Anwendung}
\abhakgruppe{Einheit 2 · Rate und Bestand als Paar}
\abhak{11}{Rate und Bestand als Paar}
\abhak{12}{Zeitraum der Abnahme eines Bestands über Nullstellen und Vorzeichen der Änderungsrate berechnen}
\abhak{13}{Rate und Bestand als Paar – Fehler finden, Begründen, Darstellungswechsel, Anwendung}
\abhakgruppe{Einheit 3 · Am Ratengraphen}
\abhak{14}{Am Ratengraphen}
\abhak{15}{Am Ratengraphen – Fehler finden, Begründen, Darstellungswechsel, Anwendung}
\end{abhakseite}
